
\chapter{Measurement}



\section{Distance}
Distance describes movement and ability range.

\vspace{0.4em}

Distance is measured in inches, with 1" roughly corresponding to two paces.

\vspace{0.4em}

For maximum flexibility,  
it is recommended that distance be measured dynamically using a tape measure.
These rules assume players will use this method of determining 
movement and range.



\section{Duration}
Durations that are determined with a \Roll{d10} may be measured in 
Turns or Actions (see Chapter 3: Combat), 
depending on the context of the effect.

\vspace{0.4em}

Each Turn is measured on the order of about a dozen seconds,
and an Action is roughly one-third of a Turn.



\section{Creature Size}
A character or creature's size is determined by their race or species.

\vspace{0.4em}

Character and creature sizes are approximate in a wide general range 
and considered \textit{Similarly-sized} when they are in the same overall class.

\subsection{Creature Sizes}

\begin{Table}{l X}
  \TableHeader{Size} & \TableHeader{Examples} \\
  Tiny     & Bug, Rodent, Pixie     \\
  Small    & Cat, Dog, Gnome        \\
  Normal   & Dwarf, Human, Wolf     \\
  Large    & Horse, Ogre, Elephant  \\
  Enormous & Wyvern, Giant, House   \\
  Colossal & Whale, Kraken, Dragon  \\
  Titanic  & Leviathan, Kaiju       \\
\end{Table}



\section{Weapon Size}
Weapon size is determined by how many hands the weapon takes to wield.

\vspace{0.4em}

This standard also applies to handheld items.



\section{Item Size \& Volume}
Item size and volume are measured in the context of inventory storage 
and bottles of liquid.


\subsection{Storage Sizes}
An item's \textit{Storage Class} is an approximation 
that takes into account the item's physical properties:
\begin{itemize}
  \item Size \& Dimensions
  \item Mass/Weight
  \item Accessibility (meaning, 
    \textit{"how easy is the item to find while rummaging through pockets?"})
\end{itemize}

\begin{Table}{l X}
  \TableHeader{Size} & \TableHeader{Description} \\
  Tiny     & Extremely small item which is generally 
    part of a collection (e.g. a bead, coin) \\
  Small    & Handheld item (e.g. a flask, bag of 33 coins), 
    or \textbf{33 Tiny items} \\
  Medium   & Larger handheld item 
    (e.g. a book, ingot of metal, bag of 100 coins), 
    or \textbf{3 Small items} \\
  Large    & Bulky adventuring item 
    (e.g. a bedroll, tent, armor, coil of rope), 
    or \textbf{3 Medium items} \\
\end{Table}



\subsection{Bottle Sizes}
Liquids, including potions are measured in \textit{Charges}. 
The potency of potions and poisons are described by the number of charges 
that can be contained in a Flask.

\vspace{0.4em}

One charge is consumed each time the bottle is used.

\vspace{0.4em}

Bottles can only contain one kind of liquid at a time.

\begin{Table}{l l X}
  \TableHeader{Size} & \TableHeader{Storage Class} 
    & \TableHeader{Description} \\
  Flask   & Small  & Smallest bottle size \\
  Bottle  & Medium & Equivalent to 3 flasks \\
  Jug     & Large  & Equivalent to 3 bottles \\
\end{Table}



